\begin{lemma}[Situation 7.3, Claim \(j-1\): nonempty case]
Fix a rung-local minimum hierarchy counterexample \(C\), set

\[
 j=\noderef{HierarchyPathNumber}(C.\mathit{hierarchy}),\qquad
 \ell=\noderef{HierarchyFullTailLength}(C),
\]

and assume \(2\le j<\ell\).  Assume every vertex has a missing color
under every proper exactly-one-uncolored replacement coloring.  If
\(\noderef{HierarchyEligibleCommonMissingIndices}\) is nonempty, then
there are an actual \noderef{HierarchyReplacementPresentation} \(P\)
and a color \(\alpha\) missing under \(P.\mathit{coloring}\) both at the
terminal vertex and somewhere in the exact prefix \(T(y_{j-2})\), such
that one of the following holds.  Writing \(H=P.H\), \(q\) for its last
stored index, \(D=D_{n,q}\),
\(\Gamma^q=\bigcup_{\eta\in D}P.\mathit{good}.\Gamma^q_\eta\), and
\(M_0\) for the colors missing on
\(P.\mathit{replacementBase}.\mathit{core}-T_n\), either

\[
                       \alpha\notin\Gamma^q\cup M_0,          \tag{A}
\]

or there is an \(\eta\in D\) such that
\(\alpha\in P.\mathit{good}.\Gamma^q_\eta\) and the first vertex missing
\(\eta\) occurs no later than \(y_{j-2}\).
\end{lemma}
\begin{proof}
Let \(I\) be the finite set
\noderef{HierarchyEligibleCommonMissingIndices}.  Its assumed
nonemptiness gives a greatest member \(i\).  Thus \(i\in I\), and the
already normalized bound
\noderef{HierarchySituationThreeEligibleIndexBound} gives
\[
                         1\leq i<j.                           \tag{1}
\]
Unfolding membership in \(I\), choose its actual
\noderef{HierarchyReplacementPresentation} \(P\) and its color
\(\alpha\).  Thus, for \(\mu=P.\mathit{coloring}\),
\[
 \alpha\in\noderef{MissingColors}(G,\mu,y_i)\cap
              \noderef{MissingColors}(G,\mu,y_\ell).          \tag{2}
\]
Write \(\mathcal E=P.\mathit{corresponding}.\mathit{replacementETT}\),
\(H=P.\mathit{counterexample}.\mathit{hierarchy}\), and let \(q\) be the
last stored index.  Let \(n\) be the last rung index of \(\mathcal E\),
put \(T_n=\mathcal E.\mathit{series}.\mathit{data}.
                 \mathit{tree}(n.\mathit{castSucc})\),
write \(D_n=\noderef{TashkinovStoredColors}
(\mathcal E.\mathit{series}.\mathit{data},n)\), and put
\[
 A=\noderef{HierarchyStarredStage}(H,q),\qquad
 B=\noderef{HierarchyFullTailPrefix}(C,j-2),\qquad
 B^+=\noderef{HierarchyFullTailPrefix}(C,j-1),
\]
\[
 D=\noderef{HierarchyStoredColorsAt}(H,q),\qquad
 g=P.\mathit{counterexample}.\mathit{good},\qquad
 \Gamma^q=\bigcup_{\eta\in D}g.\mathit{gamma}(q,\eta),
\]
\[
 M_0=\bigcup_{v\in V(P.\mathit{replacementBase}.\mathit{core})
                         \setminus V(T_n)}
                 \noderef{MissingColors}(G,\mu,v).            \tag{3}
\]
The underlying \noderef{SameHierarchyStages} clause of the
\noderef{SameGoodHierarchy} field of \(P\) identifies every ordered
stage, divider position, and exact tail prefix with the original data.
Its remaining clauses identify \(D\) and every paper-defined Gamma pair.
In particular, all occurrences of \(A,B,B^+\), \(j\), and \(\ell\) below
refer to the same individual vertices and parallel-edge identities in
the original and replacement presentations.

Apply \noderef{HierarchyCounterexamplePenultimateElementary} to the six
successive fields of the
\noderef{HierarchyMinimumCounterexampleContext}
\(P.\mathit{minimumContext}\): properness and exact one-uncoloredness,
the lower-rung package, fewer-divider induction, same-rung lexicographic
minimality, and the \(2n+11\) estimate.  We obtain the derived fact
\[
 \noderef{ColoredVertexSetElementary}
   (G,\mu,\noderef{TreeSequenceVertexSet}(T(y_{\ell-1}))).    \tag{4}
\]
The exact-prefix and strictness fields of
\noderef{TashkinovHierarchy} put \(A\prec T(y_i)\preccurlyeq
T(y_{\ell-1})\).  Since \(\alpha\) is missing at \(y_i\), elementarity
(4) says that it is missing at no vertex of \(A\).  At the last stored
index, unfolding \noderef{HierarchyStoredColorsAt} and
\noderef{TashkinovStoredColors} gives the literal equality
\[
 D=D_n\setminus\noderef{TreeMissingColors}(\mu,A).            \tag{5}
\]
The same definitions give a universal separation that will be used in
both color-choice branches.  Fix arbitrary \(\eta\in D\) and
\(\xi\in D\cap g.\mathit{gamma}(q,\eta)\).  Equation (5) says that
\(\xi\) is not missing on \(A\).  On the other hand, the avoidance field
of \noderef{GoodHierarchy}, read through
\noderef{HierarchyStageMissingColors}, puts \(\xi\) in the last
stage-missing set.  If \(q>0\), this is the missing set on the last
starred stage \(A\).  If \(q=0\), it is the missing set on \(T_n\); in
the PE case the parallel clause of \noderef{HierarchyBase} and its
\noderef{OrderedTreeUnion} make \(T_n\) an ordered initial segment of
\(A\), while in the non-PE case the nonparallel clause gives \(A=T_n\).
Thus in every case \(\xi\) is missing on \(A\), a contradiction.  Since
\(\eta\) and \(\xi\) were arbitrary,
\[
 D\cap g.\mathit{gamma}(q,\eta)=\varnothing
                         \quad\hbox{for every }\eta\in D.     \tag{5a}
\]

The same stage containment, together with the fact that \(\alpha\) is
not missing on \(A\), excludes \(\alpha\) from every last Gamma pair.
Also, the core and stage clauses of \noderef{HierarchyBase} and
\noderef{TashkinovHierarchy} put every vertex of the base core outside
\(T_n\) in \(A\).  Therefore
\[
                    \alpha\notin\Gamma^q\cup M_0.            \tag{6}
\]
If \(i\leq j-2\), the same witness \(y_i\) belongs to \(B\), and
\(P,\alpha\), and (6) prove the conclusion.  Henceforth (1) leaves only
\[
                              i=j-1.                          \tag{7}
\]
In particular, elementarity (4) and the least-candidate definition
\noderef{HierarchyFirstMissingPosition} show that the first
\(\alpha\)-missing position is exactly the last vertex of \(B^+\).

We next prove the paper's finite color choice.  Apply \noderef{NineN} to
\(\mathcal E\), its replacement base, \(H\), its goodness witness, (4),
and the \(2n+11\) field of \(P.\mathit{minimumContext}\), at the index
\(j-2\).  The inequalities \(2\leq j<\ell\) give the required index
bounds.  After unfolding its available and tail-used sets, let
\[
 R=\noderef{TreeMissingColors}(\mu,B)\setminus
       \bigl(M_0\cup\mu\langle B-A\rangle
                    \cup(\Gamma^q\cup D)\bigr).              \tag{8}
\]
If \(|R|\geq5\), the exact one-entry extension \(B\prec B^+\) adds one
individual edge and hence uses at most one color.  Choose
\(\beta\in R\) different from that edge color.  Then
\[
 \beta\in\noderef{TreeMissingColors}(\mu,B),\qquad
 \beta\notin M_0\cup\mu\langle B^+-A\rangle
                     \cup(\Gamma^q\cup D).                  \tag{9a}
\]

Suppose instead that \(|R|\leq4\).  The second conclusion of
\noderef{NineN} supplies an injection
\(h\mapsto\eta_h\) from seven indices into \(D\), with every
\(\eta_h\) missing on \(B\), such that
\[
 S_h=g.\mathit{gamma}(q,\eta_h)\cup\{\eta_h\}
       \quad\hbox{is disjoint from }\mu\langle B-A\rangle.   \tag{9b}
\]
We verify, rather than assume, that the seven sets \(S_h\) are pairwise
disjoint.  For \(h\ne h'\), injectivity gives
\(\eta_h\ne\eta_{h'}\), and the disjoint-pairs field of
\noderef{GoodHierarchy} makes their Gamma pairs disjoint.  The universal
separation (5a), applied once to each owner, gives both cross-disjointness
relations and keeps each stored singleton out of both Gamma pairs; the
two singleton stored colors are distinct by injectivity.  Hence
\(S_h\cap S_{h'}=\varnothing\).

The one new edge of \(B^+-B\) has only one color, so its color lies in at
most one of these seven pairwise disjoint \(S_h\).  Choose \(m\) for
which \(S_m\) avoids that edge.  Then \(S_m\) avoids the whole exact
interval \(B^+-A\).  The two-colors field of
\noderef{GoodHierarchy} says that
\(g.\mathit{gamma}(q,\eta_m)\) has cardinality two;
choose a color \(\beta\) in this pair.  Avoidance makes \(\beta\)
stage-missing and hence missing on \(B\), while the universal separation
(5a) gives
\(\beta\notin D\).  We have
\[
 \eta_m\in D\cap\noderef{TreeMissingColors}(\mu,B),\quad
 \beta\in g.\mathit{gamma}(q,\eta_m),\quad
 S_m\cap\mu\langle B^+-A\rangle=\varnothing.                \tag{9c}
\]
Because \(\eta_m\) has a missing witness in \(B\), unfolding
\noderef{HierarchyFirstMissingPosition} gives
\[
       (\operatorname{First}_{\mu}(\eta_m)).\mathit{val}
                                                     \leq |B|.\tag{11}
\]
Thus the chosen \(\beta\), in either (9a) or (9c), is missing on \(B\),
does not belong to \(D\), and is unused on the exact interval
\(B^+-A\).  In (9a) it lies outside \(\Gamma^q\cup M_0\); in (9c) it
lies in the Gamma pair indexed by the stored color \(\eta_m\) satisfying
(11).  Since \(\alpha\) is missing at \(y_{j-1}\), \(\beta\) is missing
at an earlier vertex of \(B\), and both vertices lie in the elementary
prefix (4), we also have \(\alpha\ne\beta\).

We record two pointwise facts used in the three recolorings.  First,
\noderef{ChangeLemma}, with the lower-rung, fewer-divider, minimality,
penultimate-elementarity, and large-set fields of the current minimum
context, identifies the component through the two relevant first missing
witnesses.  The component through \(y_\ell\) is disjoint from that
component in the sense of \noderef{TwoColorComponent}.  Indeed, equality
would put the two distinct first witnesses and \(y_\ell\), which are
distinct by vertex injectivity in \noderef{TreeSequence}, in one
component.  By
\noderef{MissingColorIncidenceEndpoint}, all three vertex-nodes have
incidence degree at most one in
\noderef{TwoColorIncidenceGraph}.  The component is spanned by a simple
path by \noderef{TwoColorComponentSpanningPath}, whereas
\noderef{TwoColorIncidenceDegreeAtMostTwo} allows only the two endpoints
of that path to have degree at most one.  This contradiction and
\noderef{TwoColorComponentsDisjoint} give the required disjointness; the
same endpoint lemma gives the possibly singleton
\noderef{TwoColorPathComponent} through \(y_\ell\).

Second, let a current proper coloring \(\psi\) have elementary
\(T(y_{\ell-1})\), and let \(\psi'\) be obtained by switching such a
disjoint terminal component \(Q\) in colors \(c,d\).  Then
\[
 \noderef{MissingColors}(G,\psi',v)=
 \noderef{MissingColors}(G,\psi,v)
             \quad(v\in V(T(y_{\ell-1}))).                  \tag{12}
\]
If \(\operatorname{inl}v\notin Q\), connectedness in
\noderef{TwoColorIncidenceGraph} keeps every incident edge-node out of
\(Q\), so \noderef{TwoColorComponentSwap} fixes each incident value.  If
\(\operatorname{inl}v\in Q\), component disjointness and elementarity
say that neither \(c\) nor \(d\) is missing at \(v\).  Expanding
\noderef{PresentColors}, choose the incident edge of each color.
\noderef{IsProperEdgeColoring} makes each identity unique, connectedness
puts both edge-nodes
in \(Q\), and the component swap exchanges their values while fixing all
other incident colors.  This proves (12) even when \(Q\) meets the prefix
at internal vertices.

We now construct, rather than posit, the exact augmentation certificates
needed in every Cases 1/2 application below.  Let \(R\) be the current
replacement presentation, \(\psi=R.\mathit{coloring}\), and suppose a
terminal \((c,d)\)-switch gives \(\psi'\), fixes every displayed edge
from \(A\) through \(X\), and satisfies (12).  Apply
\noderef{StableChangeCasesOneTwo} to this exact presentation and switch.
Its conclusion supplies
\[
 \noderef{HierarchyWeaklyStable}(R.H,q,\psi'),\qquad
 \noderef{ModColoring}(R.\mathcal E,\psi').                  \tag{13}
\]
Inspect the actual terminal certificate
\[
 R.\mathcal E.\mathit{series}.\mathit{certified}
       (\operatorname{Fin.last}R.\mathcal E.\mathit{rounds})
\]
from \noderef{TashkinovSeries} and
\noderef{TashkinovSeriesStep}.  In PE its
\noderef{TashkinovClosure} contains the old-color
\noderef{TashkinovAugmentationSequence} from the terminal ladder tree
\(T_n^R\).  In SE/RE the constructor supplies an individual connecting
edge \(f_n^R\), the literal singleton equation \(F_n^R=\{f_n^R\}\), an
exact \noderef{TreeSequenceAdjoinEdge} seed \(T_n^R+f_n^R\), and the
route from that seed.  The below-final field of
\noderef{ExtendedTashkinovTree} places \(X\preccurlyeq T\) below its
output.

In every positive-divider case, hierarchy nesting puts the applicable
source below \(A\); for non-PE this uses
\noderef{HierarchyNonparallelFirstEdge} to retain the exact seed.  At
divider zero in PE, the parallel clause of \noderef{HierarchyBase} and
its \noderef{OrderedTreeUnion} put \(T_n^R\preccurlyeq A\).  At divider
zero in SE/RE, the nonparallel clause gives \(A=T_n^R\), and the exact
one-edge seed is retained as the source.  Apply
\noderef{TashkinovAugmentationSequenceRestrict} to the certified route,
first at \(T\), then at \(X\), and, when the source precedes \(A\), at
\(A\).  Induction on the reflexive-transitive definition of
\noderef{TashkinovAugmentationSequence}, using the unit length increase
in each \noderef{TashkinovAugmentation}, splits the route at its unique
members of lengths \(|A|\) and \(|X|\).  Equality of same-length exact
\noderef{TreeSequenceInitialSegment}s of \(T\) identifies these members
with \(A\) and \(X\).  We
therefore obtain the old-color source-to-\(A\), \(A\)-to-\(X\), and
suffix \(X\longrightarrow T\) routes.  In divider-zero SE/RE, the first
piece is omitted and the exact seed-to-\(X\) route is retained.

If the divider count is positive, unfold
\noderef{HierarchyWeaklyStable} in (13).  Its positive branch contains
the \noderef{TreeColoringInvariant} on \(A\), fixing every individual
displayed edge and every vertex missing set there.  The root value is
also fixed because both colorings leave exactly \(e_0\) uncolored.
Thus \noderef{TashkinovAugmentationSequenceColoringTransport} transports
the source-to-\(A\) route.  At divider zero in PE, (13) instead unfolds
to \noderef{CombinedStableColoring} on \(T_n^R\) and its closure.
Apply \noderef{CombinedStableClosurePreserved}; it preserves that exact
closure and fixes all displayed edges and vertex missing sets on both
constituents.  The left-segment and new-entry clauses of the
\noderef{OrderedTreeUnion} cover, one individual identity at a time,
every edge and vertex of \(A\).  This supplies all premises needed to
transport the \(T_n^R\)-to-\(A\) route.  At divider zero in SE/RE there is
no route below the retained seed; its individual edge is fixed by the
non-PE \noderef{HierarchyWeaklyStable} relation in (13), because the
actual \noderef{TashkinovSeriesStep} makes that same edge incident with
\(T_n^R\).

Now use \noderef{TashkinovAugmentationSequenceColoringTransport} on the
\(A\)-to-\(X\) piece, or on the whole seed-to-\(X\) piece in the last
case.  The switch-specific exact edge equality supplies its first
pointwise hypothesis, and (12), restricted to the exact penultimate
prefix of \(X\), supplies its missing-set hypothesis.  Concatenation in
the reflexive-transitive relation gives
\[
 \begin{cases}
 T_n^R\longrightarrow X,&\text{PE},\\
 T_n^R+f_n^R\longrightarrow X,&\text{SE/RE},
 \end{cases}                                                  \tag{14}
\]
under \(\psi'\), retaining the singleton edge equation in the second
line.  The disjoint witness component leaves both affected colors
missing on \(X\) under \(\psi'\).  Hence
\noderef{TashkinovAugmentationSequenceComponentSwap} transports the old
suffix \(X\longrightarrow T\) to \(\psi'\).  Concatenating this suffix
with (14) gives the complete actual terminal route.  These are exactly
the \(hrouteX\), \(hsuffix\), and split-missing certificates used below;
no equality of all suffix edge colors is asserted.

The output of \noderef{StableChangeCasesOneTwoSplitTail} is relative to
the current \(\mathcal E_R\).  Compose its
\noderef{CorrespondingExtendedTashkinovTree} with
\(R.\mathit{corresponding}\), field by field: compose rung equalities,
dependent ladder presentations, ladder-tree equalities, raw stored-set
equalities, exact connecting-edge sets, extension kinds, and the
positive-stage stable-coloring relations.  The last compositions use
\noderef{sc2}; equality of the controlled sets follows from the raw data,
terminal-ladder missing sets, and the definition
\noderef{TashkinovStoredColors}.  Compose the two
\noderef{SameGoodHierarchy} witnesses through their
\noderef{SameHierarchyStages}, stored-set, and Gamma-pair fields.

If one affected color is thereby missing both at its retained prefix
witness and at \(y_\ell\), the two distinct vertices show that the
returned terminal tree is not elementary, so they define a
\noderef{HierarchyCounterexample}.  Rebuild its
\noderef{HierarchyMinimumCounterexampleContext} field by field.
Properness is \noderef{TwoColorComponentSwapProper}; expanding
\noderef{TwoColorComponentSwap} shows that \(\mathsf{none}\) is fixed,
so \noderef{IsOneEdgeUncolored} says that exactly \(e_0\) remains
uncolored.  The rung equality in correspondence
transports the full lower-rung package.  Same-stage equality transports
the fewer-divider assertion.  The path number, divider count, and
terminal excess are unchanged, so the two
\noderef{HierarchyCounterexampleMeasure}s are equal and the old
\noderef{HierarchyCounterexampleLess}-minimality field applies.  Finally,
correspondence and the terminal-ladder missing-set equality transport the
\(2n+11\) estimate.  This constructs an actual replacement presentation
in the original family, not merely a recolored tree.

We first handle the case \(\alpha\notin D\).  Since \(\alpha\) is not
missing on \(A\), (5) implies \(\alpha\notin D_n\).  The first missing
position of \(\beta\) is at most \(|B|\), whereas that of \(\alpha\) is
the last vertex of \(B^+\).  Apply \noderef{ChangeLemma}, with \(\beta\)
in its earlier color slot, to identify the component through those two
first witnesses.  By the three-endpoint argument above, it is disjoint
from the \((\beta,\alpha)\)-component \(Q\) through \(y_\ell\).  Switch
\(Q\), obtaining \(\sigma\).  The switch makes \(\beta\) missing at
\(y_\ell\) and preserves its old missing witness in \(B\); (12) holds.

If \(\beta\) is not missing on the last hierarchy stage, then (9c) is
impossible because every Gamma color is stage-missing.  Thus (9a) holds.
Its exclusion of \(M_0\), together with the base-core description, says
that \(\beta\) is missing nowhere on \(A\); the first \(\beta\)-witness
therefore occurs strictly after every vertex of \(A\).  Both affected
colors lie outside \(D\), and (5) now gives
\(\alpha,\beta\notin D_n\).  The first \(\beta\)-witness precedes the
first \(\alpha\)-witness, and (9a) excludes \(\beta\) from every exact
edge between them.  These are precisely the ordered external-position,
unstored-color, and unused-interval hypotheses of
\noderef{StableChangeCaseThree}, with colors ordered as
\((\beta,\alpha)\).  That theorem returns an actual corresponding
same-rung maximum ETT on \(T\), a rebuilt base, the same good hierarchy,
and hence, by the fieldwise composition and minimum-context construction
above, a replacement presentation \(P'\).

It remains in the unstored-\(\alpha\) case that \(\beta\) is
stage-missing.  Apply \noderef{StableChangeCasesOneTwo} with \(\beta\) in
the first color slot.  At positive divider count its case hypothesis is
the stage-missing assertion itself.  At divider zero it is either
\(\beta\)-missingness on \(T_n\), or \(\beta\)-missingness on \(A\) with
both affected colors outside \(D_n\); in (9a), exclusion from \(M_0\)
reduces missingness on the parallel core to missingness on \(T_n\), and
in (9c) Gamma avoidance already puts \(\beta\) in the stage-zero missing
set.  Thus every PE, positive-divider SE/RE, and divider-zero SE/RE case
is covered by the actual terminal constructor.

Take \(X=B^+\).  The color \(\alpha\) is absent from the exact interval
\(B^+-A\) by \noderef{HierarchyTailMissingColorExclusion}; its possible
divider-zero seed exception is eliminated by \(\alpha\notin D_n\).
The color \(\beta\) is absent there by (9a) or (9c).  Since the swap
changes only these colors, every exact displayed edge from \(A\) through
\(X\) is fixed.  Equation (12) supplies pointwise missing-set equality,
and both colors remain missing on \(X\) at their separated prefix
witnesses.  The construction (13)--(14), applied to the actual terminal
\noderef{TashkinovSeriesStep} of \(P\), gives the PE route from \(T_n\),
the positive-divider SE/RE route from its certified seed through \(A\),
or the divider-zero SE/RE route from the retained exact seed.

We verify the two quantified hierarchy premises.  Fix arbitrary
\(\eta\in D\).  Both affected colors are outside \(D\), so
\(\eta\ne\alpha,\beta\).  The component swap fixes the missing predicate
for \(\eta\) at every vertex.  Unfold
\noderef{HierarchyFirstMissingPosition}: if its candidate set is
nonempty, the two identical predicates give the same finite candidate
set and least member; if it is empty, both definitions return the same
last-position fallback.  Hence
\[
 \operatorname{First}_{\sigma}(\eta)
                         =\operatorname{First}_{\mu}(\eta).  \tag{15}
\]
For the Gamma cutoff, (6) puts \(\alpha\) in no last Gamma pair.  In
(9a), \(\beta\) is also in none, so the implication is vacuous.  In
(9c), if \(\beta\in\Gamma^q_\eta\), pairwise disjointness in
\noderef{GoodHierarchy} and
\(\beta\in\Gamma^q_{\eta_m}\) force \(\eta=\eta_m\); (11) then gives
\(\operatorname{First}_{\mu}(\eta).\mathit{val}\leq|B|<|B^+|\).
Thus (15) and the cutoff discharge the universal \(hstoredFirst\) and
\(hgammaCutoff\) premises of
\noderef{StableChangeCasesOneTwoSplitTail}.  Apply that theorem with the
certificates just constructed and package its output as \(P'\) by the
fieldwise construction above.

In both the external and stage-missing subcases, \(P'.\mathit{coloring}\)
has \(\beta\) missing at \(y_\ell\) and on \(B\).  Equation (12), the
same-stage data, and the unchanged base core show that its set \(M_0\)
equals the old one.  Hence (9a) gives alternative (A).  In (9c), the
same-good data retain the pair \(\Gamma^q_{\eta_m}\), while (15) retains
the first-position bound (11); this gives the second alternative in the
statement.  Thus \(P',\beta\) proves the claim whenever
\(\alpha\notin D\).

Finally suppose \(\alpha=\eta_t\in D\).  The two-colors field of
\noderef{GoodHierarchy} supplies two distinct colors
\(\gamma_1,\gamma_2\) in \(\Gamma^q_{\eta_t}\).  Its avoidance field,
read via \noderef{HierarchySegmentColors}, and the fact that the first
\(\eta_t\)-missing position is \(y_{j-1}\), give
\[
 \gamma_1,\gamma_2\notin\mu\langle B^+-A\rangle.             \tag{16}
\]
Equation (5a) gives \(\gamma_1,\gamma_2\notin D\), and in particular
neither can equal \(\eta_t\).
The theorem \noderef{HierarchyTailMissingColorExclusion} excludes
\(\eta_t\) from the same interval, with one possible exception: at
divider zero in SE/RE, the first displayed edge may be the individual
stored connecting seed.

Apply \noderef{ChangeLemma} to \((\gamma_1,\eta_t)\).  It identifies the
component through the stage-missing first witness of \(\gamma_1\) with
the component through \(y_{j-1}\).  Separate this component from the
terminal component \(Q_1\) by the three-endpoint argument and switch
\(Q_1\), obtaining \(\mu_1\).  Apply
\noderef{StableChangeCasesOneTwo} with \(\gamma_1\) in its first slot.
For positive divider count \(\gamma_1\) is missing on the last ordinary
stage; for divider zero it is missing on \(T_n\).  Thus its case premise
holds in PE, SE, and RE.

We now verify all split-tail inputs for this first switch, with
\(X=B^+\).  Equation (16) fixes every possible \(\gamma_1\)-edge on the
exact interval.  Tail exclusion fixes every \(\eta_t\)-edge except the
divider-zero SE/RE seed.  In that exceptional branch, unfold the
non-PE \noderef{HierarchyWeaklyStable} conclusion of the just-applied
stable-change theorem.  It is a \noderef{StableColoring} relation on
\(T_n=A\).
The actual SE or RE constructor of \noderef{TashkinovSeriesStep} gives
the singleton connecting edge and its endpoint in \(T_n\); the incident
edge clause of that stable relation fixes this same individual seed.
Hence every displayed edge from \(A\) through \(B^+\) is fixed in all
three terminal kinds.  The pointwise argument (12) gives missing-set
equality on the penultimate prefix, and the separated witness component
leaves both \(\gamma_1\) and \(\eta_t\) missing on \(B^+\).

For \(hstoredFirst\), fix \(\eta\in D\).  If
\(\eta\ne\eta_t\), then \(\eta\) differs from both affected colors, so
its missing predicate is globally fixed; equality of candidate sets and
the common terminal fallback give equality of first positions exactly as
in (15).  If \(\eta=\eta_t\), equation (12) identifies its missing
predicate at every vertex through \(T(y_{\ell-1})\).  Its old first
witness \(y_{j-1}\) lies in that prefix, so the candidate sets have the
same least member there; no vertex after the prefix can change that least
member.  Thus the first position is again unchanged.  This treats every
stored color, including the totalized fallback in the unaffected case.

For \(hgammaCutoff\), the universal separation (5a) says that the stored
color \(\eta_t\) belongs to no Gamma pair.  If
\(\gamma_1\in\Gamma^q_\eta\), then its simultaneous membership in
\(\Gamma^q_{\eta_t}\) and pairwise disjointness force
\(\eta=\eta_t\).  Its first missing position is \(y_{j-1}\), so
\[
       \operatorname{First}_{\mu}(\eta).\mathit{val}=|B^+|. \tag{17}
\]
The route construction (13)--(14), now instantiated with the actual
terminal certificate of \(P\), supplies the PE, positive-divider SE/RE,
and divider-zero SE/RE route alternatives.  Its transported suffix is
obtained specifically from
\noderef{TashkinovAugmentationSequenceComponentSwap}.  Apply
\noderef{StableChangeCasesOneTwoSplitTail} using the just-proved edge,
route, missing-set, stored-position, and cutoff premises.  Compose its
correspondence and same-good data with those of \(P\), and rebuild all
minimum-context fields as above.  This gives an actual original-relative
replacement presentation \(P_1\), not only the coloring \(\mu_1\).
Under \(\mu_1\), the color \(\gamma_1\) is missing both at its retained
stage witness and at \(y_\ell\), and every missing set on the
penultimate prefix agrees with that under \(\mu\).

Apply \noderef{HierarchyCounterexamplePenultimateElementary} afresh to
the six fields of \(P_1.\mathit{minimumContext}\).  Thus
\(T(y_{\ell-1})\) is elementary under \(\mu_1\); this is not borrowed
from the old presentation.  The first switch fixes all displayed edges
through \(B^+\), preserves all missing sets there pointwise, and its
same-good data retain \(D\) and all Gamma pairs.  Consequently (9a), or
(9c) together with (11), remains valid under \(\mu_1\).

First note that \(\gamma_1\ne\beta\).  In (9a),
\(\beta\notin\Gamma^q\), whereas
\(\gamma_1\in\Gamma^q_{\eta_t}\).  In (9c), the stored colors
\(\eta_m\) and \(\eta_t=\alpha\) have distinct missing witnesses in
\(B\) and at \(y_{j-1}\), respectively; elementarity gives
\(\eta_m\ne\eta_t\), and the disjoint-pairs field of
\noderef{GoodHierarchy} separates their Gamma pairs.

In the \(P_1\) presentation, name the two first missing witnesses
\[
 w_\gamma=T.\mathit{vertex}
       (\operatorname{First}_{\mu_1}(\gamma_1)),\qquad
 w_\beta=T.\mathit{vertex}
       (\operatorname{First}_{\mu_1}(\beta)).
\]
The color \(\gamma_1\) retains a missing witness in \(A\), so the
least-candidate definition \noderef{HierarchyFirstMissingPosition} and
the exact initial-segment relation put \(w_\gamma\) in \(A\).  Likewise,
(9a) or (9c) supplies a \(\beta\)-missing witness in \(B\), so
\(w_\beta\) lies in \(B\).  Hence both lie in \(B\).  Apply
\noderef{ChangeLemma} to \((\gamma_1,\beta)\).  It identifies the
components through \(w_\gamma\) and \(w_\beta\); call their common
component \(Q_B\).  We prove that \(Q_B\) is disjoint from the component
\(Q_2\) through \(y_\ell\), splitting on equality of the two named
witnesses.

If \(w_\gamma=w_\beta=w\), then both \(\gamma_1\) and \(\beta\) are
missing at \(w\).  Unfolding \noderef{MissingColors}, no edge incident
with \(w\) has either affected color, by the incidence-witness definition
\noderef{PresentColors}.  By
\noderef{TwoColorIncidenceGraph}, the vertex-node \(\operatorname{inl}w\)
therefore has no adjacent edge-node, and the bipartite definition gives
no other possible neighbor.  Thus it is isolated.  Unfolding
\noderef{TwoColorComponent}, every walk from this node has length zero,
so
\[
                         Q_B=\{\operatorname{inl}w\}.
\]
The exact-prefix definition puts \(w\) in the strict prefix \(B\), while
the terminal vertex \(y_\ell\) is outside \(B\); vertex injectivity in
\noderef{TreeSequence} gives \(w\ne y_\ell\).  Hence \(Q_B\ne Q_2\),
because \(\operatorname{inl}y_\ell\in Q_2\), and
\noderef{TwoColorComponentsDisjoint} makes them disjoint.

If \(w_\gamma\ne w_\beta\), then
\(w_\gamma,w_\beta,y_\ell\) are three distinct vertices: the first two
lie in \(B\), whereas the terminal vertex is outside \(B\).  If
\(Q_B=Q_2\), all three vertex-nodes lie in that component.  Each is
missing one of the two affected colors, so
\noderef{MissingColorIncidenceEndpoint} gives incidence degree at most
one at all three.  But \noderef{TwoColorComponentSpanningPath} supplies
a simple path spanning the component, and
\noderef{TwoColorIncidenceDegreeAtMostTwo} is the ambient two-color
degree bound.  Every internal node of the spanning path has its two
distinct path neighbors, so a node of degree at most one must be one of
the path's two endpoints.  Three distinct such nodes are impossible.
Thus \(Q_B\ne Q_2\), and
\noderef{TwoColorComponentsDisjoint} again makes the two components
disjoint.

Only now switch \(Q_2\), obtaining \(\mu_2\), and apply
\noderef{StableChangeCasesOneTwo} with \(\gamma_1\) first.  It is still
stage-missing at its separated witness, so the case premise holds at
positive divider count and at divider zero.

Again take \(X=B^+\).  The first switch and (16) show that
\(\gamma_1\) is unused on \(B^+-A\) under \(\mu_1\); equations (9a) or
(9c) do the same for \(\beta\).  Thus the second switch fixes every exact
edge through \(X\).  The pointwise internal-vertex argument (12), now
using the elementarity just rederived for \(P_1\), gives
\[
 \noderef{MissingColors}(G,\mu_2,v)=
 \noderef{MissingColors}(G,\mu_1,v)
             \quad(v\in V(T(y_{\ell-1}))).                  \tag{18}
\]
Both affected colors remain missing on \(X\) at their separated prefix
witnesses.

Every member of the current stored set \(D\) differs from
\(\gamma_1\) and \(\beta\): the universal separation (5a) excludes
\(\gamma_1\), and it also excludes \(\beta\) in (9c), while (9a)
excludes \(\beta\) directly.  Hence the component swap fixes its missing
predicate globally.  Unfolding
\noderef{HierarchyFirstMissingPosition} gives identical candidate sets
and least elements, or the identical last-position fallback.  This proves
\(hstoredFirst\) for the second switch.  For its Gamma cutoff,
\(\gamma_1\in\Gamma^q_{\eta_t}\), and disjoint pairs show that any Gamma
pair containing \(\gamma_1\) has index \(\eta_t\); its first position is
still bounded by \(|B^+|\) by (17) and the first-switch stored-position
equality.  In case (9a), \(\beta\) lies in no pair.  In case (9c), any
pair containing \(\beta\) is indexed by \(\eta_m\), whose first position
is at most \(|B|<|B^+|\) by (11) and the same first-switch equality.
This proves the cutoff for every stored index and both affected colors.

Finally inspect the actual terminal
\noderef{TashkinovSeriesStep} of
\(P_1.\mathit{corresponding}.\mathit{replacementETT}\), not that of the
original presentation.  Repeat the certified construction
(13)--(14): positive-divider \noderef{TreeColoringInvariant} transports
the route below \(A\); divider-zero PE uses
\noderef{CombinedStableClosurePreserved} and the exact
\noderef{OrderedTreeUnion}; divider-zero SE/RE retains its singleton
seed; pointwise equality (18) transports the finite prefix; and
\noderef{TashkinovAugmentationSequenceComponentSwap} transports the
actual old suffix.  Apply
\noderef{StableChangeCasesOneTwoSplitTail} with these route certificates,
the stored-position equality, and the two cutoff bounds.  Compose the
returned correspondence first with \(P_1.\mathit{corresponding}\) and
then with the original one, and likewise compose the
\noderef{SameGoodHierarchy} data.  Rebuild properness,
one-uncoloredness, the lower package, fewer-divider induction,
lexicographic minimality, and the large estimate field by field.  The
result is an actual presentation \(P_2\) in the original replacement
family.

The second switch makes \(\beta\) missing at \(y_\ell\) and preserves
its witness in \(B\).  Equations (12) and (18), together with equality of
the base core, identify the final \(M_0\) with the original one.  Thus
(9a) yields alternative (A).  In (9c), composed same-goodness retains
\(\beta\in\Gamma^q_{\eta_m}\), and the two stored-position equalities
retain (11).  Hence \(P_2,\beta\) satisfy the second alternative.  This
proves the claim also when \(\alpha\in D\).
\end{proof}
